\documentclass[11pt,a4paper]{article}
\usepackage[T1]{fontenc}
\usepackage[utf8]{inputenc}
\usepackage[czech]{babel}
\usepackage{amsmath,amssymb,amsthm}
\usepackage[margin=2.5cm]{geometry}
\newtheorem{theorem}{Věta}
\newtheorem{lemma}[theorem]{Lemma}
\newcommand{\E}{\mathbb E}
\newcommand{\HH}{\mathsf H}
\title{Přenos ocasu mezi větvemi se zachováním logaritmického účtu}
\author{}
\date{19. září 2026}
\begin{document}
\maketitle
\noindent\textbf{Rozsah.} Dokazujeme přesný kladný přesun mezi konečnou
hlavou a ocasem různých dětí a účet platný v libovolné hloubce při
libovolných přípustných přesunech. Existence politiky, která takto udrží
všechny budoucí fany nezáporné, zde dokázána není. Ze samotného účtu
nevyvozujeme logaritmickou horní mez pro původní úlohu.

\section*{1. Stav s proměnnou amplitudou ocasu}
Pišme $P=(8,16,16,16,8)/64$, $Q=(7,20,10,20,7)/64$,
$q_d=Q^{*d}$ a $W_a(j)=a/((a+j-1)(a+j))$ pro $j\ge1$.
Pro $0\le s\le4d$ definujme konečně vážený, ale nekonečně podporovaný ocas
\[
 C_{d,s}=\sum_{y=0}^{4d}\frac{2q_d(y)}{8d+2-y},
             \delta_{8d-s}*W_{8d+2-y},\qquad
 T_d=|C_{d,s}|=\E\frac{2}{8d+2-Y_d},\quad Y_d\sim q_d.
\]
Obecný orbitální stav nyní má tvar
\begin{equation}\label{state}
 B_h=A_h+\kappa_h C_{d,s(h)},\qquad
 A_h\ge0,\quad \operatorname{supp}A_h\subseteq[0,8d-s(h)],
 \quad |A_h|=1-\kappa_hT_d,
\end{equation}
kde $\kappa_h\ge0$. Kořen je $A_\varnothing=0$, $\kappa_\varnothing=1$,
$C_{0,0}=W_2=H_2$.

Rozbalení jednotlivých Waringových fází dává
\[
 Q*C_{d,s}=D_{d,s}+\delta_x*C_{d+1,s+x}
 \qquad\text{pro každé }x\in\{0,\ldots,4\}.
\]
Pravý ocas po posunutí o $x$ nezávisí na $x$.
Míra $D_{d,s}\ge0$ je podporována na $8d-s+1,\ldots,8d-s+8$,
má hmotnost $T_d-T_{d+1}$ a explicitně pro $1\le v\le8$
\[
 D_{d,s}(8d-s+v)=
 2\sum_y q_d(y)\sum_{z=0}^{\min(4,v-1)}
 \frac{Q_z}{(8d+2-y+v-z-1)(8d+2-y+v-z)}.
\]

\begin{lemma}[Úplná rovnice rozšířené orbity]\label{fan}
Stavy \eqref{state} splňují exactní fan právě tehdy, když
\begin{equation}\label{system}
 \boxed{
 \sum_xP_x\delta_x*A_{hx}=Q*A_h+\kappa_hD_{d,s(h)},
 \qquad \sum_xP_x\kappa_{hx}=\kappa_h.}
\end{equation}
Při podmínkách \eqref{state} je to kladný causal fan s čerstvým vstupem $Q$.
\end{lemma}
\begin{proof}
Ocas za společným řezem $8(d+1)-s(h)$ má na obou stranách stejný
nenulový profil. Rovnost jeho amplitud je druhá rovnice
\eqref{system}. Po jeho odečtení zbude první rovnice. Naopak jejich
sečtení dává celý fan koeficient po koeficientu. Zdroj lze nejprve
vylosovat jako součet aktuální rezervy a nezávislého $Q$; příslušná
nezáporná transportní matice poté určí výstup.
\end{proof}

\section*{2. Explicitní kladný přesun hlava--ocas}
Nechť již máme jeden takový fan a jeho konečnou transportní matici
$\pi_x(u)=P_xA_{hx}(u-x)$. Označme $t=T_{d+1}$ a vezměme různé řádky
$i,j$ a konečný sloupec $u\ge\max(i,j)$. Pro
\[
 0\le\varepsilon\le
 \min\{\pi_i(u),\ P_j\kappa_{hj}t\}
\]
proveďme
\begin{align*}
 A_{hi}'&=A_{hi}-\frac{\varepsilon}{P_i}\delta_{u-i},&
 \kappa_{hi}'&=\kappa_{hi}+\frac{\varepsilon}{P_it},\\
 A_{hj}'&=A_{hj}+\frac{\varepsilon}{P_j}\delta_{u-j},&
 \kappa_{hj}'&=\kappa_{hj}-\frac{\varepsilon}{P_jt}.
\end{align*}
Ostatní děti ponecháme.

\begin{theorem}[Přesun mezi konečnou zásobou a ocasem]\label{exchange}
Tento předpis zachovává nezápornost, normalizaci všech dětí a celý
exactní fan. Zvyšuje rezervu v dítěti $i$ ve stochastickém pořadí a
snižuje ji v dítěti $j$; přijímající ocas leží nad konečným sloupcem.
\end{theorem}
\begin{proof}
Omezení na $\varepsilon$ přesně zajišťují nezápornost odečítaných měr.
V každém dítěti se změna hmotnosti hlavy ruší se změnou hmotnosti ocasu.
Ve zdrojovém sloupci $u$ se změny po násobení $P_i,P_j$ vyruší.
Také změna $\sum_xP_x\kappa_{hx}$ je nulová.
Lemma \ref{fan} dokazuje exactnost. Podpora ocasu začíná až za
společným konečným řezem, tedy nad $u$.
\end{proof}

Například v již známém kvantilovém fanu z $H_2$ zvolíme
\[
 i=1,\quad j=0,\quad u=2,\quad\varepsilon=\frac1{128}.
\]
Jeho ocas má hmotnost $T_1=10327/40320$. Nové amplitudy jsou
\[
 \boxed{\kappa_0=\frac{7807}{10327},\qquad
 \kappa_1=\frac{11587}{10327},\qquad
 \kappa_2=\kappa_3=\kappa_4=1.}
\]
V konečných hlavách se odečte $1/32$ v rezervě $1$ dítěte $1$ a přičte
$1/16$ v rezervě $2$ dítěte $0$. Proto se veličina
$27A_1(0)+7A_1(1)$ zlepší o $7/32$, zatímco tatáž veličina dítěte $0$
se nezmění. Změna jeho vyšší části samozřejmě může ovlivnit pozdější
pokračovatelnost; ta se tímto výpočtem neprokazuje.

\section*{3. Účet nezávislý na rozdělení amplitud}
Nechť libovolný kladný exactní strom tvaru \eqref{state} existuje do
hloubky $n$. Pro $|h|=d$ označme $p(h)=\prod_iP_{h_i}$ a
\[
 M_d=\sum_{|h|=d}p(h)\sum_j jA_h(j),\qquad
 S_d=\sum_{|h|=d}p(h)\kappa_hs(h),\qquad
 V_d=2\E\HH_{8d+2-Y_d},
\]
kde $\HH_k=\sum_{j=1}^k1/j$.

\begin{theorem}[Invariantní logaritmický účet]\label{account}
Pro každou hloubku $0\le d\le n$ platí přesně
\begin{equation}\label{exactaccount}
 \boxed{
 \sum_{|h|=d}p(h)\kappa_h=1,\qquad
 M_d-T_dS_d=V_d-(8d+1)T_d-2.}
\end{equation}
Zejména, pro $d\ge1$, rovnoměrně přes všechny takové stromy,
\begin{equation}\label{log}
 \boxed{M_d=2\log d+O(1).}
\end{equation}
Amplitudy $\kappa_h$ nemusejí mít společnou horní mez přes historie.
\end{theorem}
\begin{proof}
První rovnost v \eqref{exactaccount} je indukce z druhé rovnice
\eqref{system}. Proto také $0\le S_d\le4d$.

Pro pevná celá $a\ge1,r\ge0$ platí
\[
 \mathbb E\min(r+W_a,N)
 =a\HH_N+r+1-a\HH_a+o(1).
\]
Pro náhodnou veličinu $J_h$ se zákonem \eqref{state} tedy
\[
 \mathbb E\min(J_h,N)=2\kappa_h\HH_N+c_h+o(1),
\quad c_h=\sum_jjA_h(j)+\kappa_h(8d-s(h)+1)T_d-\kappa_h V_d.
\]
V jednom exactním fanu se vedoucí členy ruší díky
$\sum_xP_x\kappa_{hx}=\kappa_h$. Omezený posun náhodné veličiny o $y$
přidává ke konstantnímu členu právě $y$, a $\mathbb EP=\mathbb EQ=2$.
Proto $c_h=\sum_xP_xc_{hx}$.
V kořeni je $c_\varnothing=-2$, takže vážený součet $c_h$ v libovolné
pevné hloubce je $-2$. Dosazení a první rovnost v
\eqref{exactaccount} dokazují druhou rovnost.

Nakonec $4d+2\le8d+2-Y_d\le8d+2$ a
$T_d\le2/(4d+2)$. Z $0\le S_d\le4d$ plyne například
\[
 2\HH_{4d+2}-6\le M_d\le2\HH_{8d+2}-2.
\]
To dokazuje \eqref{log}. V každé pevné hloubce je pouze konečně mnoho
uzlů a fází, takže jsme nikde nevyměňovali nekonečný součet přes
historie s limitou ani nepoužili uniformní těsnost.
\end{proof}

\section*{4. Co tato konstrukce uzavírá a co ne}
Přesun z věty \ref{exchange} je skutečná změna kladného fanu.
Věta \ref{account} dokazuje, že libovolné přípustné rozdělování amplitud
v této rozšířené orbitě zachová logaritmický účet hlav. Účet tedy
nevyžaduje společný ocas ani společné $\kappa$ všech dětí.

Není však dokázáno, že tyto přesuny vždy dovolí pokračovat až do další
hloubky. Proto z těchto vět zatím neplyne existence svědků pro všechna
$n$. Kdyby takoví svědci s kořenem $H_2$ existovali, jejich kořenovou
rezervu lze pro horizont $n$ omezit na $4n$: v původním protokolu
uchováme původní skrytou hodnotu, ale skutečnou rezervu inicializujeme
jejím minimem s $4n$. Při původní hodnotě větší než $4n$ tato skutečná
rezerva vystačí na libovolných $n$ výstupů; jinak protokol beze změny
kopírujeme. Tím by vznikla horní mez
\[
 \mathcal B_n\le\mathbb E\min(H_2,4n)
 =2(\HH_{4n+1}-1)=2\log n+O(1).
\]
Tento poslední závěr zůstává podmíněný existencí kladných svědků pro
neomezené horizonty.
\end{document}
